% Appendix E — N6 ternary X2MH6 2/2 falsification; cation-VEC insufficient.
\section{Appendix E --- The $X_2MH_6$ ternary funnel: cation-VEC is
necessary but insufficient}
\label{app:E}

The N5 wall (Appendix~\ref{app:D}) motivated a ternary funnel: stuffed-cation
octahedral hydrides $X_2MH_6$ (Fm$\bar{3}$m, K$_2$PtCl$_6$ prototype) that
might decouple coupling from stability. This appendix reports the
\emph{pre-registered} falsification of that route at ambient pressure.

\paragraph{Structure and the octahedral molecular orbitals.}
The Fm$\bar{3}$m K$_2$PtCl$_6$ prototype places the transition metal $M$ at
the $4a$ Wyckoff site $(0,0,0)$, the cations $X$ at $8c$
$(\frac14,\frac14,\frac14)$, and the six hydrogens at $24e$ $(x,0,0)$ with
$x\approx0.24$ --- a 9-atom primitive cell ($1M + 2X + 6H$). The six H~$1s$
orbitals form the symmetry-adapted ligand combinations $a_{1g}+e_g+t_{1u}$,
which hybridise with the metal $s$/$p$/$d$ manifold to produce bonding,
non-bonding, and antibonding ($\sigma^\ast$) bands. The $\sigma^\ast$ band is
H--H/H--M antibonding and therefore steep in the density of states; its
partial filling is the engine of the proposed superconductivity.

\paragraph{The cation-VEC rule.}
The valence-electron count of the octahedral unit is the closed-form
\begin{equation}
\mathrm{VEC} = 2\,v(X) + v(M) + 6,
\qquad
n_{\sigma^\ast} = \mathrm{VEC} - 18,
\end{equation}
where $v(\cdot)$ is the group valence, the $+6$ is the six H electrons, and
$n_{\sigma^\ast}$ is the $\sigma^\ast$ filling above the octahedral
18-electron closed shell. The model predicts a sweet spot at
VEC$\in\{19,20\}$ ($n_{\sigma^\ast}=1$--2): at VEC$=18$ the shell is closed
($N(E_F)\to0$, $T_c\to0$); at VEC$\ge21$ the $\sigma^\ast$ over-fills, the
M--H bond weakens, and a soft mode collapses the lattice. The
$\{19,20\}$ window is the intersection of ``enough electrons'' and ``stable
lattice.'' Both \ce{Mg2IrH6} ($2\cdot2+9+6=19$) and \ce{Li2CuH6}
($2\cdot1+11+6=19$) land exactly on it --- so the rule, as a heuristic, is
self-consistent. The question this appendix answers is whether VEC$=19$ is
\emph{sufficient}.

\begin{table}[h]
\centering
\small
\begin{tabular}{ccl}
\toprule
\textbf{VEC} & \textbf{$n_{\sigma^\ast}$} & \textbf{state} \\
\midrule
18 & 0 & closed shell, $N(E_F)\downarrow$, $T_c\approx0$ \\
19 & 1 & $\sigma^\ast$ quarter-filled, $N(E_F)\uparrow$, strong e--ph (sweet spot) \\
20 & 2 & $\sigma^\ast$ half-filled, optimal $\lambda$ (sweet spot) \\
21 & 3 & over-filled, M--H weakened $\to$ soft mode, unstable \\
22 & 4 & severe over-fill $\to$ lattice decomposition \\
\bottomrule
\end{tabular}
\caption{Octahedral $\sigma^\ast$ band-filling model. The campaign-tested
VEC$=19$ candidates sit in the predicted sweet spot, yet are dynamically
unstable at ambient pressure (Tab.~\ref{tab:ternary}) --- the rule is a
necessary filter, not a stability guarantee.}
\label{tab:vec}
\end{table}

\paragraph{Why VEC alone is insufficient: the cation-radius axis.}
\ce{Ca2IrH6} and \ce{Sr2IrH6} also have VEC$=19$ ($2\cdot2+9+6$), yet are
expected to have $T_c\approx0$. The second axis is the cation \emph{ionic
radius}: small, electronegative cations ($\mathrm{Li}^+$ $0.76$\,\AA,
$\mathrm{Mg}^{2+}$ $0.72$\,\AA) chemically pre-compress the \ce{H6} cage,
stiffening the H phonons; large cations ($\mathrm{Ca}^{2+}$ $1.00$\,\AA,
$\mathrm{Sr}^{2+}$ $1.18$\,\AA) expand the lattice, soften the H phonons,
lower $\omega_{\log}$, and let cation $d$-states intrude near $E_F$, diluting
the hydrogen-dominated coupling. The honest closed-form statement of this is fenced (it is a
DFT/wet-lab mechanism question, not an identity):
\begin{lstlisting}[basicstyle=\ttfamily\footnotesize,frame=single,xleftmargin=0pt]
$ hexa verify --fence "X2MH6 Tc peaks at VEC 19-20 (sigma* partial fill
  maximizes N(EF)); Mg->Ca/Sr keeps VEC=19 but expands lattice, softens H
  phonons and intrudes cation d-states, killing Tc"
  tier = WHITE SPECULATION-FENCED  (mechanism, verification N/A by design)
\end{lstlisting}

\paragraph{Pre-registered falsifiers.}
We registered the pass/fail threshold before harvesting either result:
a candidate is dynamically stable (and the sweet spot ``sufficient'') only
if every $q$-point clears $\texttt{min\_freq} > -\SI{50}{\per\centi\meter}$
with at most one hard-imaginary mode. Tab.~\ref{tab:ternary} is the verdict.

\begin{table}[h]
\centering
\small
\begin{tabular}{lcccl}
\toprule
\textbf{Candidate} & \textbf{VEC} & \textbf{ambient stable?} &
  \textbf{min freq (\si{\per\centi\meter})} & \textbf{verdict} \\
\midrule
\ce{Mg2IrH6} & 19 & \textbf{no} (48\% hard-imag) & $-2235$ & \tierRed\ FALSIFIED \\
\ce{Li2CuH6} & 19 & \textbf{no} ($8$ hard-imag modes) & $-944.92$ & \tierRed\ FALSIFIED \\
\bottomrule
\end{tabular}
\caption{The $X_2MH_6$ cation-VEC$=19$ sweet spot, falsified 2/2 at ambient
pressure. \ce{Mg2IrH6}: 5/13 partial $q$-verdict already violates the
threshold (PR \#247). \ce{Li2CuH6}: the second $q$-point alone locks 8
hard-imaginary modes (PR \#275). The K$_2$PtCl$_6$ Fm$\bar{3}$m prototype is
ambient-unstable for both.}
\label{tab:ternary}
\end{table}

\paragraph{The falsifier statements (verbatim, F-N6 ledger).}
\begin{quote}\small
\textbf{F-N6-1.} If \ce{Mg2IrH6} Fm$\bar{3}$m at $P=0$ satisfies
$\texttt{min\_freq} > -\SI{50}{\per\centi\meter}$ on all $q$ with
hard-imag count $\le 1$, the VEC$=19$ sweet spot is \emph{sufficient};
otherwise it is only \emph{necessary}.\\
\textbf{Result:} \tierRed\ PASSED 2026-05-26 --- threshold clearly violated
($-\SI{2235}{\per\centi\meter}$, 48\% hard-imag). Sufficiency falsified.

\medskip
\textbf{F-N6-2.} The same threshold for \ce{Li2CuH6}; failure
family-falsifies the K$_2$PtCl$_6$ prototype at ambient.\\
\textbf{Result:} \tierRed\ PASSED 2026-05-27 --- $q_2$ alone: 8 hard-imag,
$-\SI{944.92}{\per\centi\meter}$. $X_2MH_6$ Fm$\bar{3}$m 2/2 falsified.
\end{quote}

\paragraph{The cation-decoupling escape (open route).}
The N5 wall (Appendix~\ref{app:D}) and the falsification above point to the
same escape: decouple the coupling numerator $\eta = N(E_F)\langle
I^2\rangle$ from the stability denominator $\langle\omega^2\rangle$.

\begin{figure}[h]
\centering
\begin{tikzpicture}[font=\footnotesize,node distance=6mm]
\node[draw,rounded corners,fill=red!8,text width=4.4cm,align=center] (b) {%
  \textbf{binary} $H_3X$\\[2pt]
  $\eta$ and $\langle\omega^2\rangle$ both tied to one knob\\
  $\Rightarrow$ soft for $\lambda\uparrow$ but unstable};
\node[draw,rounded corners,fill=teal!10,text width=4.8cm,align=center,
      right=14mm of b] (t) {%
  \textbf{ternary} $X_2MH_6$ / clathrate\\[2pt]
  cation doping: $N(E_F)$ preserved ($\eta\uparrow$)\\
  cation pre-compression: cage stiffens ($\langle\omega^2\rangle\uparrow$)\\
  $\Rightarrow$ two independent knobs};
\draw[-{Latex},thick] (b) -- (t) node[midway,above,font=\tiny]{decouple};
\end{tikzpicture}
\caption{The cation-decoupling hypothesis. In binary hydrides a single knob
($\langle\omega^2\rangle$) controls both coupling and stability; a
stuffed-cation phase can in principle move them independently. \ce{CaH6}
($m=+0.5$, Apx.~\ref{app:D}) is the existence proof; the ambient
$X_2MH_6$ octahedral instances tested here ($m<0$, unstable) are not.}
\label{fig:decouple}
\end{figure}

\noindent The expansion matrix derived from the VEC rule (small cation,
group-9/11 metal, VEC$\in\{19,20\}$) nominates \ce{Mg2RhH6}, \ce{Mg2CoH6},
\ce{Al2MnH6} (VEC$=19$) and \ce{Mg2NiH6}, \ce{Mg2PtH6} (VEC$=20$) as
not-yet-falsified candidates, and pre-excludes \ce{Ca2IrH6}, \ce{Sr2IrH6},
\ce{K2CuH6}, and the VEC$\ge21$ over-filled set to save DFT budget. These
are \emph{rule-predicted only} (no DFT yet) and carry no $T_c$ claim.

\paragraph{Reading.}
The cation-VEC rule survives as a necessary filter (it correctly places both
candidates on the sweet spot and explains why \ce{Mg}$\to$\ce{Ca}/\ce{Sr}
substitution kills $T_c$ despite preserving VEC --- the larger ionic radius
expands the lattice and softens the H phonons). But VEC$=19$ does \emph{not}
guarantee dynamical stability: the octahedral $\sigma^\ast$ filling that
maximises $N(E_F)$ also drives the soft modes that collapse the lattice. This
is the N5 wall reappearing in the ternary family. The honest consequence is
that the ambient $X_2MH_6$ octahedral route, as tested, is closed; the
clathrate $MXH_8$ route (\ce{LaBeH8} anchor) remains open and is the
campaign's recommended next probe (\S\ref{sec:conclusion}).
